% AGU 2026 abstract — Opti O2 / Beaver Creek WA dissolved-oxygen classification
% Derived from the contents and conclusions of slides.py (the presentation deck).
%
% AGU limits: Title <= 300 characters, Abstract text <= 2000 characters.
% Plain-text counts of the content below: Title = 197 chars, Abstract = 1908 chars.
% (The counts refer to the rendered prose, not the LaTeX source markup.)
%
% Compile: pdflatex agu2026_abstract.tex

\documentclass[11pt]{article}
\usepackage[utf8]{inputenc}
\usepackage[T1]{fontenc}
\usepackage{amsmath}
\usepackage[margin=1in]{geometry}
\usepackage{parskip}
\usepackage{lineno}

\title{Tidal Hiccups and Oxic Sighs: Machine-Learning classification of
oxygenation events in a transitioning coastal wetland}

\author{Tanmay Pani}
\date{AGU Annual Meeting 2026}

\begin{document}
\maketitle

\begin{abstract}
%	\linenumbers
Coastal wetlands are converting former freshwater floodplains into tidally
influenced marsh, reshaping the subsurface dissolved-oxygen (DO) regime that
governs biogeochemical function. At Beaver Creek, Washington, a near-anoxic
floodplain well (DO levels below detection $\sim$92\% of the time), punctuated by
episodic oxygenation events. These events fall into two physically distinct
classes: abrupt, asymmetric ``hot moments'' driven by tidal/saline intrusion, and
slower, symmetric ``oxic pulses'' driven by freshwater rise from precipitation/flooding. 
We formalize separating these regimes as a supervised
classification problem using six years of co-located 5-minute DO, hydrology
(water level, salinity, temperature), and weather measurements.

We compare two modeling tracks. The first uses tabular classification models on domain-engineered event-shape features encoding morphology of the DO vs time data, change in salinity and water-level,
antecedent-precipitation, etc. The second directly uses raw multivariate time-series data either as inputs to a Convolutional Neural Network (CNN), or through automatic encoding into \(\mathcal{O}(10^4)\) features capturing correlations over various time scales, fed into tabular models. Classification performance is reported through metrics like Confusion Matrix, F1-score, Balanced Accuracy, etc.

Furthermore, through feature importance attribution, we have are isolating and identifying the various drivers for these DO events.
\end{abstract}

\end{document}
